\documentclass[UTF8,12pt,a4paper,fontset=fandol]{ctexbook}


% 支持从项目根目录或当前笔记目录编译。
\makeatletter
\def\input@path{{../}{../../}{../../../}{../../../../}}
\makeatother
\usepackage{researchnotes}


\title{高等代数（谢启鸿）笔记}
\author{}
\date{\today}


\begin{document}


\maketitle
\tableofcontents
\clearpage


\chapter{行列式}

\section{二阶行列式}

\section{三阶行列式}

\section{\texorpdfstring{$n$}{n} 阶行列式}

\section{行列式的展开和转置}

\section{行列式的计算}

\section{行列式的等价定义}

\section{Laplace 定理}


\chapter{矩阵}

\section{矩阵的概念}

\section{矩阵的运算}

\section{方阵的逆阵}

\section{矩阵的初等变换与初等矩阵}

\section{矩阵乘积的行列式与初等变换法求逆阵}

\section{分块矩阵}

\section{Cauchy-Binet 公式}


\chapter{线性空间}

\section{数域}

\section{行向量和列向量}

\section{线性空间}

\section{向量的线性关系}

\section{向量组的秩}

\section{矩阵的秩}

\section{坐标向量}

\section{基变换与过渡矩阵}

\section{子空间}

\section{线性方程组的解}


\chapter{线性映射}

\section{线性映射的概念}

\section{线性映射的运算}

\section{线性映射与矩阵}

\section{线性映射的像与核}

\section{不变子空间}


\chapter{多项式}

\section{一元多项式代数}

\section{整除}

\section{最大公因式}

\section{因式分解}

\section{多项式函数}

\section{复系数多项式}

\section{实系数多项式和有理系数多项式}

\section{多元多项式}

\section{对称多项式}

\section{结式和判别式}


\chapter{特征值}

\section{特征值和特征向量}

\section{对角化}

\section{极小多项式与 Cayley-Hamilton 定理}

\section{特征值的估计}


\chapter{相似标准型}

\section{多项式矩阵}

\section{矩阵的法式}

\section{不变因子}

\section{有理标准型}

\section{初等因子}

\section{Jordan 标准型}

\section{Jordan 标准型的进一步讨论和应用}

\section{矩阵函数}


\chapter{二次型}

\section{二次型的化简与矩阵的合同}

\section{二次型的化简}

\section{惯性定理}

\section{正定型与正定矩阵}

\section{Hermite 型}


\chapter{内积空间}

\section{内积空间的概念}

\section{内积的表示和正交基}

\section{伴随}

\section{内积空间的同构、正交变换和酉变换}

\section{自伴随算子}

\section{复正规算子}

\section{实正规矩阵}

\section{谱分解与极分解}

\section{奇异值分解}

\section{最小二乘解}


\chapter{双线性型}

\section{对偶空间}

\section{双线性型}

\section{纯量积}

\section{交错型与辛几何}

\section{对称型与正交几何}


\end{document}
